%2multibyte Version: 5.50.0.2960 CodePage: 1252

\documentclass[english,12pt,a4paper]{article}
%%%%%%%%%%%%%%%%%%%%%%%%%%%%%%%%%%%%%%%%%%%%%%%%%%%%%%%%%%%%%%%%%%%%%%%%%%%%%%%%%%%%%%%%%%%%%%%%%%%%%%%%%%%%%%%%%%%%%%%%%%%%%%%%%%%%%%%%%%%%%%%%%%%%%%%%%%%%%%%%%%%%%%%%%%%%%%%%%%%%%%%%%%%%%%%%%%%%%%%%%%%%%%%%%%%%%%%%%%%%%%%%%%%%%%%%%%%%%%%%%%%%%%%%%%%%
\usepackage{setspace,amsmath,amsfonts,amssymb,amsthm}
\usepackage{multirow}
\usepackage{graphicx,subfig}
\usepackage{epstopdf}
\usepackage[margin=1.0in]{geometry}
\usepackage{setspace}
\usepackage[capposition=top]{floatrow}
\usepackage{indentfirst}
\usepackage{booktabs,dcolumn}
\usepackage{hyperref}

\setcounter{MaxMatrixCols}{10}
%TCIDATA{OutputFilter=LATEX.DLL}
%TCIDATA{Version=5.50.0.2960}
%TCIDATA{Codepage=1252}
%TCIDATA{<META NAME="SaveForMode" CONTENT="1">}
%TCIDATA{BibliographyScheme=Manual}
%TCIDATA{LastRevised=Wednesday, May 13, 2026 14:13:03}
%TCIDATA{<META NAME="GraphicsSave" CONTENT="32">}

\onehalfspace
\input{tcilatex}
\begin{document}

\title{\textbf{Sovereign Default: The Role of Expectations}\thanks{
We thank Patrick Kehoe for many useful conversations. We also thank the
associate editor and two referees, as well as Fernando Alvarez, Manuel
Amador, Cristina Arellano, V.V. Chari, In-Koo Cho, Jonathan Eaton, Carlo
Galli, Veronica Guerrieri, Jonathan Heathcote, Andy Neumeyer, Fabrizio
Perri, Martin Uribe, and Vivian Yue for very useful comments. Nicolini and
Teles would like to acknowledge the support of FCT as well as the ADEMU
project, \textquotedblleft A Dynamic Economic and Monetary
Union,\textquotedblright\ funded by the European Union's Horizon 2020
Program under grant agreement N${\ {}^{\circ }}$ 49396. The views expressed
herein are those of the authors and not necessarily those of the Federal
Reserve Bank of Minneapolis, the Federal Reserve System, Banco de Portugal,
the European System of Central Banks, or the Inter-American Development
Bank. E-mail addresses: joaoay@iadb.org, gaston.m.navarro@frb.gov,
juanpa@minneapolisfed.org, pteles@ucp.pt.}}
\author{Jo\~{a}o Ayres \\
%EndAName
{\normalsize {Inter-American Development Bank}}\\
[1ex] Gaston Navarro \\
[-0.3ex] {\normalsize {Federal Reserve Board} }\\
[1ex] Juan Pablo Nicolini \\
[-0.3ex] {\normalsize {Federal Reserve Bank of Minneapolis and Universidad
Di Tella} }\\
[1ex] Pedro Teles \\
[-0.3ex] {\normalsize {Banco de Portugal, Catolica Lisbon School of Business
\& Economics, CEPR}}}
\date{September 2017}
\maketitle

\begin{abstract}
In the standard model of sovereign default, as in Aguiar and Gopinath (2006)
or Arellano (2008), default is driven by fundamentals alone. There is no
independent role for expectations. We show that small variations of that
model are consistent with multiple interest rate equilibria, similar to the
ones found in Calvo (1988). For distributions of output that are commonly
used in the literature, the high interest rate equilibria have properties
that make them fragile. Once output is drawn from a distribution with both
good and bad times, however, it is possible to have robust high interest
rate equilibria. \smallskip

\noindent \textit{Keywords: } Sovereign default; multiple equilibria; good
and bad times.

\noindent \textit{JEL Codes: }E44, F34.

\medskip \bigskip
\end{abstract}

\section{ A two-period model of IMF\ Catalytic role}

We study a two-period endowment economy populated by a representative agent
that can borrow from a continuum of risk-neutral foreign creditors. The
utility function of the representative agent is given by 
\begin{equation*}
\ln c_{1}+E_{1}\ln c_{2}
\end{equation*}%
Income at time 1 is assumed to be equal to $e$ which we assume to be small
enough so the agent wants to save.

In the second period, the agent receives endowment $y$, which follows 
\begin{equation*}
y_{2}=%
\begin{cases}
y^{l}, & \text{ with probability }p \\ 
y^{h}, & \text{ with probability }\left( 1-p\right)%
\end{cases}%
\end{equation*}

In period one, the representative agent can borrow $b$ in a noncontingent
bond in international financial markets.

The risk-neutral gross international interest rate is $R^{\ast }$. \ The
interest rate charged to the country is given by $R.$

We also allow for an international organization (LOLR) that is willing to
lend an amount $h$, at an interest rate $R^{I}.$

In period two, after observing the realization of the endowment, the
borrower decides whether to:

\begin{enumerate}
\item Pay the debt in full. In this case, consumption is given by%
\begin{equation*}
y_{2}-bR-hR^{I}
\end{equation*}

\item Default only on international lenders. In this case output is given by 
$y^{d}$ which is lower than the expected value of $y.$ In this case,
consumption will be%
\begin{equation*}
y^{d}-hR^{I}
\end{equation*}

\item Default on the LOLR is given by a fixed cost $\kappa .$ Thus, \ if
there is default on both lenders and the financial institution, output is
equal to consumption and is given by 
\begin{equation*}
y^{d}-\kappa 
\end{equation*}
\end{enumerate}

Thus, to the extent that the LOLR lends an amount $h$ such that%
\begin{equation*}
hR^{I}\leq \kappa ,
\end{equation*}%
the country will never default on the LOLR.

Note that a similar environmnet is one in which if there is defaault on
private bonds, the endowment is $y^{d}$, while if it defaults on the IMF it
is $y^{d}-\kappa $. Thus, the country could default on private bonds only or
default on both, since there is no extra cost of defaulting on private bonds.

In what follows, we evaluate the effect on the equilibrium set of different
values of $(h,R^{I})$ satisfying the inequality above.

Then consumption in the second period depends on the default decision as
follows:

\begin{eqnarray*}
c &=&y-hR^{I}-bR\text{ \ \ \ if pays} \\
c &=&y^{d}-hR^{I}\text{ }\ \ \ \ \ \ \ \ \ \ \text{if default }
\end{eqnarray*}%
Then, the country will default when 
\begin{equation*}
y^{d}+bR\geq y.
\end{equation*}%
Notice that if 
\begin{equation*}
y^{d}+bR>y^{2},
\end{equation*}%
the ex-ante probability of default is equal to one, so the natural debt
limit is given by 
\begin{equation*}
bR\leq y^{2}-y^{d}
\end{equation*}

Since creditors are risk neutral, in order for them to be indifferent
between lending at $R$ or at the risk-free rate $R^{\ast }$, the expected
return from lending to the borrower must be the same as $R^{\ast }$, so%
\begin{eqnarray*}
R &=&R^{\ast }\text{ \ \ \ \ \ \ \ \ \ \ \ \ \ if }bR^{\ast }\leq y^{1}-y^{d}
\\
R &=&\frac{R^{\ast }}{(1-p)}\ \ \ \text{if }y^{1}-y^{d}\leq b\frac{R^{\ast }%
}{(1-p)}\leq y^{2}-y^{d}
\end{eqnarray*}% 
or%
\begin{eqnarray*}
R &=&R^{\ast }\text{ \ \ \ \ \ \ \ \ \ \ \ if \ }b^{\ast }\leq \frac{%
y^{1}-y^{d}}{R^{\ast }} \\
R &=&\frac{R^{\ast }}{(1-p)}\ \ \ \text{if \ }\frac{y^{1}-y^{d}}{R^{\ast }}%
(1-p)\leq b\leq \frac{y^{2}-y^{d}}{R^{\ast }}(1-p)
\end{eqnarray*}
This defines a schedule $b=\left( R\right) $ .

We assume that 
\begin{equation*}
(1-p)y^{2}+py^{d}>y^{1}
\end{equation*}%
so%
\begin{equation*}
(1-p)y^{2}+py^{d}-y^{d}>y^{1}-y^{d}
\end{equation*}%
\bigskip or%
\begin{equation*}
\left( 1-p\right) \left( y^{2}-y^{d}\right) >y^{1}-y^{d}
\end{equation*}%
\bigskip This ensures the schedule is like in the Figures of the JET paper.

We focus on 2 different schedules.

Schedule 1 is given by%
\begin{eqnarray*}
R &=&R^{\ast }\text{ \ \ \ \ \ \ \ \ \ \ \ if \ }b^{\ast }\leq \frac{%
y^{1}-y^{d}}{R^{\ast }} \\
R &=&\frac{R^{\ast }}{(1-p)}\ \ \ \text{if \ }b>\frac{y^{1}-y^{d}}{R^{\ast }}
\end{eqnarray*}% 
while the other schedule is given by 
\begin{eqnarray*}
R &=&R^{\ast }\text{ \ \ \ \ \ \ \ \ \ \ \ if \ }b^{\ast }\leq \frac{%
y^{1}-y^{d}}{R^{\ast }}(1-p) \\
R &=&\frac{R^{\ast }}{(1-p)}\ \ \ \text{if \ }b>\frac{y^{2}-y^{d}}{R^{\ast }}%
(1-p)
\end{eqnarray*}

Given this schedule, the optimal choice of debt by the borrower at period 1
is to chose $h$ and $b$ to maximize utility: 
\begin{equation}
V(R(s),\kappa ,\omega )=\ln (\omega +h+b)+\beta p\max \{\ln
(y^{1}-hR^{I}-bR(s)),\ln (y^{d}-hR^{I})\}+\beta p\ln (y^{1}-hR^{I}-bR(s)).
\label{OP1}
\end{equation}% 
subject to 
\begin{equation}
bR\leq y^{2}-y^{d}  \label{MaxB}
\end{equation}% 
\begin{equation}
hR^{\ast }\leq \kappa ,  \label{MaxA}
\end{equation}

\subsection{The case in which $R^{I}=R^{\ast }$}

In considering the problem of the country at time zero, we first consider
the case in which $R^{I}=R^{\ast }.$ Thus, as investors would never lend at
rates lower than $R^{\ast },$ it is - weakly - optimal for the country to
borrow from the LOLR\ up the maximum feasible amount.

Then, we can write the problem as to chose $b$ to maximize utility: 
\begin{equation}
V(R(s),\kappa ,\omega )=\ln (\omega +\frac{\kappa }{R^{\ast }}+b)+\beta
p\max \{\ln (y^{1}-\kappa -bR(s)),\ln (y^{d}-\kappa )\}+\beta p\ln
(y^{1}-\kappa -bR(s)).  \label{OP2}
\end{equation}% 
subject to 
\begin{equation*}
bR\leq y^{2}-y^{d}
\end{equation*}% 
The lending by the LOLR\ is equivalent to increasing the endowment at time
zero $(\omega )\,$\ and reducing it at time 1. In this case it does it at
actuarially fair prices. We assume that $\omega $ is small enough relative
to $y^{1},y^{2}$ so that the country borrows from the IMF all the way to $%
\frac{\kappa }{R^{\ast }}.$

\bigskip 

In this case, it is trivial to show that an increase in $\kappa $ (weakly)
reduces $b,$ which implies that the LOLR\ can certainly play a "catalytic"
role and give incentives to the country to do endogenous austerity instead
of gambling for redemption. 

To solve this case for given $\kappa $, we proceed as follows:

1. Assume $R=R^{\ast }$ and chose $b$ to max utility ignoring constraint $%
\left( \ref{MaxB}\right) $. If the optimal solution for $b$ is lower than $%
\frac{y^{1}-y^{d}}{R^{\ast }}(1-p)$, which is the upper bound of the region
without multplicity, this is the equilibrium.

2. If not, we need to specify sunspot depenndent policies. 

3. To do so, consider first the bad sunspot. First compute utility setting $%
b=\frac{y^{1}-y^{d}}{R^{\ast }}(1-p)$ and $R=R^{\ast }.$ This is the utility
the country obtains with the bad sunspot if it does endogenous austerity.
Then, assume $R=\frac{R^{\ast }}{1-p}$ and chose the max over choice of $b$
subject to constraints $\left( \ref{MaxB}\right) $ and $\left( \ref{MaxA}%
\right) .$ This is the utility the country obtains if going beyond the jump
in rates. The equilibrium is given by the pair $(b,R)$ corresponding to the
case that delivers higher utility.

4. Consider now the good sunspot. First compute utility setting $b=\frac{%
y^{1}-y^{d}}{R^{\ast }}$ and $R=R^{\ast }.$ This is the utility the country
obtains with the good sunspot if it does endogenous austerity. Then, assume $%
R=\frac{R^{\ast }}{1-p}$ and chose the max over choice of $b$ subject to
constraints $\left( \ref{MaxB}\right) $ and $\left( \ref{MaxA}\right) .$
This is the utility the country obtains if going beyond the jump in rates.
The equilibrium is given by the pair $(b,R)$ corresponding to the case that
delivers higher utility.

\bigskip 

\subsection{The case in which $R^{I}>R^{\ast }$}

$\bigskip $

Notation: Consider the plot of the schedules offered by the lenders. Let A
be the lower bound of the set of values for teh debt such that there is
multiplicity, and let B be the upper bound of such a set.  

Imagine now the model without the IMF and assume that the equilibrium
exhibits gambling for redemption. Specifically, assume that the equilibrium
implies that with the good sunspot, the country borrows all the way to $B$
at rate $R^{\ast }$, the point at which there would be a jump in the
interest rate. And assume also that with the bad sunspot, the country
chooses $\widehat{b}\in (A,B)$ at the high rate $\frac{R^{\ast }}{1-p}.$

Now, assume that the LOLR is willing to lend to this country, at a rate  $%
R^{I}=R^{\ast }+\varepsilon ,$ for relatively low $\varepsilon ,$ an amount
up to $\widehat{b}-A.$

Consider now the following allocation: if the good sunspot happens, the
country borrows $B$ at the rate $R^{\ast },$ and it may eventually tap some
extra borrowing from the LOLR up to the amount $\frac{\kappa }{R^{\ast }}.$

In contrast, if the bad sunspot happens, the country borrows an amount $A$
from the private sector, which, for this amount, is willing to lend at the
low rate $R^{\ast }.$ Then, it borrows from the LOLR to the maximun, that is
\bigskip $\widehat{b}-A,$ at the rate $R^{I}=R^{\ast }+\varepsilon .$

The country is clearly better of, since in the good sunspot, the marginal
rate of increasing borrowing is smaller than before, while in the bad
sunspot it is borrowing the same amount at lower rates.

\bigskip 

Note that this allocation satisfies all the equilibrium conditions we
described above. However, it is clearly not an equilibrium. Private agents
have incentives to lend to the country marginal amounts above $A,$ but below 
$B\,,$ at rates $R^{\ast }<R<R^{\ast }+\varepsilon ,$ and the country would
be willing to take those loans and reduce the borrowing from the IMF\ one by
one.

Which is the \ equilibrium? 

A\ better allocation for everybody is for the private lenders to replace all
the LOLR\ lending at the rate $R,$ with  $R^{\ast }<R<R^{\ast }+\varepsilon .
$

This is still not an equilibrium yet, since private lenders are making
profits and the lender may stilll want to tap extra funds from the LOLR.....

At that allocation, private lenders want to cut down rates and they will do
so until private borrowing is equal to B.

In addition, the country may be willling to tap extra resources from the
LOLR.

Thus, the equilibrium is the same as with the good sunspot.

\bigskip 

The question is which equilibrum condition are we missing?

I think the problem is in the description of the lenders problem.

By specifying that 
\begin{equation*}
R=R^{\ast }(1-P(d=1))
\end{equation*}% 
where $P(d=1)$ is the probability of default, we are just imposing that the
return between two assets is the same. As long as those are the only two
relevant assets, that condition is necessary and sufficient. But we also
need that no other possible asset offered by the investors yields strictly
positive profits and is atractive for the country.

With the LOLR, we also need that 
\begin{equation*}
(R^{I}-R)h\leq 0
\end{equation*}%
That is, if the country is borrowing positive amounts from the LOLR, it has
to be the case that $(R^{I}-R)\leq 0.$

\bigskip 

I think we were not imposing this condition so far.

\bigskip 

\bigskip 

\end{document}
